\documentclass[handouts]{beamer}
\usepackage[orientation=portrait,size=a4, scale=2]{beamerposter}
\beamertemplatenavigationsymbolsempty

\usepackage[czech]{babel}
\usepackage{booktabs}



\title{Funkce používané v ekologii}
\author{Robert Mařík}
\usetheme{CambridgeUS}
\usecolortheme{crane}

\def\subsection#1{\par\medskip \textbf{#1}\par\medskip}
\begin{document}

% \setbeamercolor{block title}{bg=white, fg=black}
% \setbeamercolor{block body}{bg=white, fg=black}

\begin{frame}

  
  \begin{center}
    \large Funkce používané v ekologii
  \end{center}
Stavební kameny vstupující do modelů. Je vhodné je umět identifikovat a znát jejich vlastnosti.    

    \begin{columns}
    \begin{column}{0.35\textwidth}

\begin{block}{Přímá úměrnost}

        \begin{itemize}
        \item Růst konstantní rychlostí. Používá se jako výchozí předpoklad, 
        pokud nejsou důvody domnívat se, že se rychlost mění. 
        \item Matematický předpis je $$y = kx.$$ 
        \item Grafem funkce je přímka, která
         prochází počátkem 
         a bodem $[1, k]$.
        \item Je-li $k > 0$, pak funkce roste rychleji pro vyšší hodnoty $k$.
        \end{itemize}
        \includegraphics[width=\linewidth]{prima_umernost.pdf}
    
\end{block}
      \end{column}

    \begin{column}{0.6\textwidth}
\begin{block}{Kvadratická funkce}
      \begin{itemize}
      \item Růst rychlostí úměrnou velikosti populace a volnému místu v~životním prostředí. Používá se pro modelování růstu populace, pokud se předpokládá, že růst je rychlejší pro větší populace.
      \item Používá se, pokud se dynamika růstu zpomaluje a
       po překročení nosné kapacity prostředí dochází k poklesu.
      \item Matematický předpis je $$y = rx\left(1-\frac{x}{K}\right).$$
      \item Grafem funkce je parabola, která prochází počátkem a body $[0, 0]$ a 
      $[0, K]$ a je otočena vrcholem nahoru.
      \item V počátku roste rychlostí $r$.
      \end{itemize}

      \centering
    \includegraphics[width=0.55\textwidth]{kvadraticka_funkce.pdf}

\end{block}
    \end{column}
    \end{columns}

    \begin{columns}
        \begin{column}{0.48\textwidth}
\begin{block}{Mocninná funkce}
        \begin{itemize}
        \item Kombinuje jednoduchost s flexibilitou, je oblíbená 
        při formulaci alometrických vztahů.
        \item Matematický předpis je $$y = x^k.$$ Rychlost růstu se mění, pro $k>1$ rychlost růstu roste s $x$.
        \item Funkce prochází počátkem a bodem $[1,1]$, 
        pro $k>1$ má vodorovnou tečnu a pro větší $k$ je na intervalu $[0,1]$ menší a na intervalu $[1,+\infty)$ větší
        a na intervalu $[1, +\infty)$ větší. Pro $0<k<1$ má funkce v počátku svislou tečnu.
        \item Pro logaritmicky škálované osy přechází na přímou úměrnost.
        \end{itemize}

            \centering
        \includegraphics[width=.7\linewidth]{mocnina.pdf}
\end{block}
            
        \end{column}
        \begin{column}{0.46\textwidth}

\begin{block}{Lomená funkce}
    \begin{itemize}
        \item Používá se pro růst k saturované hodnotě. Například při modelování působení predátorů.
        \item Matematický předpis je $$y = \frac{x}{x+a}\quad \text{nebo}\quad y=\frac{x^k}{x^k+a^k}.$$
        \item V okolí počátku se funkce chová jako funkce $y=\frac xa$ resp. $y=\frac{x^k}{a^k}$. Pro $k>1$ má v počátku vodorovnou tečnu.
        \item Prochází bodem $[a,\frac 12]$ a zdola se blíží k vodorovné asymptotě $y=1$.
    \end{itemize}

    \centering
    \includegraphics[width=.99\linewidth]{lomena.pdf}
\end{block}
        \end{column}

        \end{columns}

\end{frame}
\end{document}


%%% Local Variables: 
%%% TeX-command-extra-options: "-shell-escape"
%%% TeX-engine: xetex
%%% End:
